% GENERATED FILE. Edit canonical lessons, then run npm run generate:books.
% canonical-source-hash: fnv1a64:66db3ca4c99d4f6f
% canonical-lessons: BN-C322-rakkhos, BN-C322-pori, BN-C322-daini, BN-C322-obhishap, BN-C322-bhut, BN-R322-review-316-322

\chapter{Ogres, Fairies and Ghosts}
\label{ch:bn-ogres-fairies-and-ghosts}
\hlchaptermodality{322}

\begin{chapteropening}
\textbf{By the end of this chapter:} \emph{I can tell of ogres, fairies, witches and curses.}

Complete the last lesson of chapter 322: I can tell of ogres, fairies, witches and curses.
\end{chapteropening}

\section[\texorpdfstring{\bn{রাক্ষস} (rākkhôs)}{রাক্ষস (rākkhôs)}]{\bn{রাক্ষস} (rākkhôs) — an ogre; a demon who eats people}

\label{lesson:BN-C322-rakkhos}



\begin{quote}
Before the new one: say the Bengali for magic, then the Bengali for to sleep.
\end{quote}

\subsection*{You'll want to know: \bn{রাক্ষস}}
\textbf{\bn{রাক্ষস}} — \emph{rākkhôs} — \textquotedblleft{}an ogre; a demon who eats people\textquotedblright{}.

\begin{quote}
\bn{রাক্ষসটা} \bn{রাতে} \bn{গ্রামে} \bn{এসে} \bn{মানুষ} \bn{খেত।}
\end{quote}

\begin{quote}
\emph{rākkhôsṭā rāte grāme ese mānush khetô.}
\end{quote}

\begin{quote}
The ogre used to come to the village at night and eat people.
\end{quote}

\begin{cousinweb}
\textbf{\bn{রাক্ষস}} (\emph{rākkhôs}) is Sanskrit \emph{rākṣasa}, a Rakṣas, a malignant demon of the old epics (Monier-Williams). Its \bn{ক্ষ} is read \emph{kkh}.
\end{cousinweb}

\begin{grammarlens}[title={What you've built}]
One more word for telling what happened.
\end{grammarlens}

\subsection*{Your turn}
\begin{itemize}
\raggedright
  \item \emph{Say it:} \emph{rākkhôs}
  \item \emph{Say it:} \emph{rākkhôs}, once more
  \item \emph{Say it:} \emph{rākkhôsṭā rāte grāme ese mānush khetô.} — The ogre used to come to the village at night and eat people.
  \item \emph{Recall:} say the Bengali for magic, then the Bengali for to sleep, then say \emph{rākkhôs} again
\end{itemize}

\begin{tcolorbox}[breakable,colback=teal!4,colframe=teal!35!black,title={Before you move on}]
What does \bn{রাক্ষস} mean? (\textquotedblleft{}An ogre; a demon who eats people\textquotedblright{}.) Say it once more.
\end{tcolorbox}
\section[\texorpdfstring{\bn{পরী} (pôrī)}{পরী (pôrī)}]{\bn{পরী} (pôrī) — a fairy}

\label{lesson:BN-C322-pori}



\begin{quote}
Before the new one: say the Bengali for to sleep, then the Bengali for an ogre.
\end{quote}

\subsection*{You'll want to know: \bn{পরী}}
\textbf{\bn{পরী}} — \emph{pôrī} — \textquotedblleft{}a fairy\textquotedblright{}.

\begin{quote}
\bn{গল্পে} \bn{একটা} \bn{পরী} \bn{শিশুটিকে} \bn{বাঁচাল।}
\end{quote}

\begin{quote}
\emph{gôlpe ekṭā pôrī shishuṭike bā̃chālô.}
\end{quote}

\begin{quote}
In the story a fairy saved the child.
\end{quote}

\begin{cousinweb}
\textbf{\bn{পরী}} (\emph{pôrī}) is usually traced to Persian \emph{parī}, 'a fairy', the winged spirit of Persian stories.
\end{cousinweb}

\begin{grammarlens}[title={What you've built}]
One more word for telling what happened.
\end{grammarlens}

\subsection*{Your turn}
\begin{itemize}
\raggedright
  \item \emph{Say it:} \emph{pôrī}
  \item \emph{Say it:} \emph{pôrī}, once more
  \item \emph{Say it:} \emph{gôlpe ekṭā pôrī shishuṭike bā̃chālô.} — In the story a fairy saved the child.
  \item \emph{Recall:} say the Bengali for to sleep, then the Bengali for an ogre, then say \emph{pôrī} again
\end{itemize}

\begin{tcolorbox}[breakable,colback=teal!4,colframe=teal!35!black,title={Before you move on}]
What does \bn{পরী} mean? (\textquotedblleft{}A fairy\textquotedblright{}.) Say it once more.
\end{tcolorbox}
\section[\texorpdfstring{\bn{ডাইনি} (ḍāini)}{ডাইনি (ḍāini)}]{\bn{ডাইনি} (ḍāini) — a witch}

\label{lesson:BN-C322-daini}



\begin{quote}
Before the new one: say the Bengali for an ogre, then the Bengali for a fairy.
\end{quote}

\subsection*{You'll want to know: \bn{ডাইনি}}
\textbf{\bn{ডাইনি}} — \emph{ḍāini} — \textquotedblleft{}a witch\textquotedblright{}.

\begin{quote}
\bn{রূপকথার} \bn{ডাইনিটা} \bn{বনের} \bn{ভিতরে} \bn{একা} \bn{থাকত।}
\end{quote}

\begin{quote}
\emph{rūpkôthār ḍāiniṭā bôner bhitôre ekā thāktô.}
\end{quote}

\begin{quote}
The witch in the fairy tale lived alone deep in the forest.
\end{quote}

\begin{cousinweb}
\textbf{\bn{ডাইনি}} (\emph{ḍāini}) comes down from Sanskrit \emph{ḍākinī}, a female spirit in the train of the goddess Kālī (CDIAL 5542, which lists Bengali \emph{ḍāin} and \emph{ḍāini}, 'a witch').
\end{cousinweb}

\begin{grammarlens}[title={What you've built}]
One more word for telling what happened.
\end{grammarlens}

\subsection*{Your turn}
\begin{itemize}
\raggedright
  \item \emph{Say it:} \emph{ḍāini}
  \item \emph{Say it:} \emph{ḍāini}, once more
  \item \emph{Say it:} \emph{rūpkôthār ḍāiniṭā bôner bhitôre ekā thāktô.} — The witch in the fairy tale lived alone deep in the forest.
  \item \emph{Recall:} say the Bengali for an ogre, then the Bengali for a fairy, then say \emph{ḍāini} again
\end{itemize}

\begin{tcolorbox}[breakable,colback=teal!4,colframe=teal!35!black,title={Before you move on}]
What does \bn{ডাইনি} mean? (\textquotedblleft{}A witch\textquotedblright{}.) Say it once more.
\end{tcolorbox}
\section[\texorpdfstring{\bn{অভিশাপ} (ôbhishāp)}{অভিশাপ (ôbhishāp)}]{\bn{অভিশাপ} (ôbhishāp) — a curse}

\label{lesson:BN-C322-obhishap}



\begin{quote}
Before the new one: say the Bengali for a fairy, then the Bengali for a witch.
\end{quote}

\subsection*{You'll want to know: \bn{অভিশাপ}}
\textbf{\bn{অভিশাপ}} — \emph{ôbhishāp} — \textquotedblleft{}a curse\textquotedblright{}.

\begin{quote}
\bn{ডাইনির} \bn{অভিশাপে} \bn{রাজকুমার} \bn{পাথর} \bn{হলেন।}
\end{quote}

\begin{quote}
\emph{ḍāinir ôbhishāpe rājkumār pāthôr hôlen.}
\end{quote}

\begin{quote}
By the witch's curse the prince was turned to stone.
\end{quote}

\begin{cousinweb}
\textbf{\bn{অভিশাপ}} (\emph{ôbhishāp}) is Sanskrit \emph{abhiśāpa}, 'a curse; a charge' (Monier-Williams). \bn{অভিশাপ} \bn{দেওয়া}, to put a curse on someone.
\end{cousinweb}

\begin{grammarlens}[title={What you've built}]
One more word for telling what happened.
\end{grammarlens}

\subsection*{Your turn}
\begin{itemize}
\raggedright
  \item \emph{Say it:} \emph{ôbhishāp}
  \item \emph{Say it:} \emph{ôbhishāp}, once more
  \item \emph{Say it:} \emph{ḍāinir ôbhishāpe rājkumār pāthôr hôlen.} — By the witch's curse the prince was turned to stone.
  \item \emph{Recall:} say the Bengali for a fairy, then the Bengali for a witch, then say \emph{ôbhishāp} again
\end{itemize}

\begin{tcolorbox}[breakable,colback=teal!4,colframe=teal!35!black,title={Before you move on}]
What does \bn{অভিশাপ} mean? (\textquotedblleft{}A curse\textquotedblright{}.) Say it once more.
\end{tcolorbox}
\section[\texorpdfstring{\bn{ভূত} (bhūt)}{ভূত (bhūt)}]{\bn{ভূত} (bhūt) — a ghost}

\label{lesson:BN-C322-bhut}



\begin{quote}
Before the new one: say the Bengali for a witch, then the Bengali for a curse.
\end{quote}

\subsection*{You'll want to know: \bn{ভূত}}
\textbf{\bn{ভূত}} — \emph{bhūt} — \textquotedblleft{}a ghost\textquotedblright{}.

\begin{quote}
\bn{লোকে} \bn{বলে} \bn{যে} \bn{পুরনো} \bn{বাড়িটায়} \bn{একটা} \bn{ভূত} \bn{থাকে।}
\end{quote}

\begin{quote}
\emph{loke bôle je purôno bāṛiṭāy ekṭā bhūt thāke.}
\end{quote}

\begin{quote}
People say a ghost lives in the old house.
\end{quote}

\begin{cousinweb}
\textbf{\bn{ভূত}} (\emph{bhūt}) is Sanskrit \emph{bhūta}, 'become, been, past' (Monier-Williams), and so a being, a spirit: one who has been.
\end{cousinweb}

\begin{grammarlens}[title={What you've built}]
One more word for telling what happened.
\end{grammarlens}

\subsection*{Your turn}
\begin{itemize}
\raggedright
  \item \emph{Say it:} \emph{bhūt}
  \item \emph{Say it:} \emph{bhūt}, once more
  \item \emph{Say it:} \emph{loke bôle je purôno bāṛiṭāy ekṭā bhūt thāke.} — People say a ghost lives in the old house.
  \item \emph{Recall:} say the Bengali for a witch, then the Bengali for a curse, then say \emph{bhūt} again
\end{itemize}

\begin{tcolorbox}[breakable,colback=teal!4,colframe=teal!35!black,title={Before you move on}]
What does \bn{ভূত} mean? (\textquotedblleft{}A ghost\textquotedblright{}.) Say it once more.
\end{tcolorbox}
\section[(dialogue)]{Looking back: from The Siege to Ogres, Fairies and Ghosts}

\label{lesson:BN-R322-review-316-322}



\begin{quote}
Say the Bengali for to surround and a rampart.
\end{quote}

\subsection*{Your turn}
Say these aloud:

\begin{itemize}
\raggedright
  \item to surround, a rampart, a moat, to hold out, to fall back
  \item a warrior, self-sacrifice, a martyr, to surrender, defeat
  \item a revolution, a movement, a struggle, to found, liberation
  \item ruins, a statue, to excavate, an inscription, currency
  \item a memorial, to construct, a tomb, a dome, architecture
  \item a legend, a fairy tale, a giant, magic, to sleep
  \item an ogre, a fairy, a witch, a curse, a ghost
\end{itemize}

\begin{tcolorbox}[breakable,colback=teal!4,colframe=teal!35!black,title={Before you move on}]
To surround? (\textbf{\bn{ঘেরা}}, \emph{gherā}.) A rampart? (\textbf{\bn{প্রাচীর}}, \emph{prāchīr}.) A moat? (\textbf{\bn{পরিখা}}, \emph{pôrikhā}.) To hold out? (\textbf{\bn{টেকা}}, \emph{ṭekā}.) To fall back? (\textbf{\bn{পিছোনো}}, \emph{pichhono}.) A warrior? (\textbf{\bn{বীর}}, \emph{bīr}.) Self-sacrifice? (\textbf{\bn{আত্মত্যাগ}}, \emph{ātmôtyāg}.) A martyr? (\textbf{\bn{শহিদ}}, \emph{shôhid}.) To surrender? (\textbf{\bn{আত্মসমর্পণ} \bn{করা}}, \emph{ātmôsômôrpôn kôrā}.) Defeat? (\textbf{\bn{পরাজয়}}, \emph{pôrājôy}.) A revolution? (\textbf{\bn{বিপ্লব}}, \emph{biplôb}.) A movement? (\textbf{\bn{আন্দোলন}}, \emph{āndolôn}.) A struggle? (\textbf{\bn{সংগ্রাম}}, \emph{sônggrām}.) To found? (\textbf{\bn{প্রতিষ্ঠা} \bn{করা}}, \emph{prôtishṭhā kôrā}.) Liberation? (\textbf{\bn{মুক্তি}}, \emph{mukti}.) Ruins? (\textbf{\bn{ধ্বংসাবশেষ}}, \emph{dhvôngsābôshesh}.) A statue? (\textbf{\bn{মূর্তি}}, \emph{mūrti}.) To excavate? (\textbf{\bn{খনন} \bn{করা}}, \emph{khônôn kôrā}.) An inscription? (\textbf{\bn{শিলালিপি}}, \emph{shilālipi}.) Currency? (\textbf{\bn{মুদ্রা}}, \emph{mudrā}.) A memorial? (\textbf{\bn{স্মৃতিস্তম্ভ}}, \emph{smritistômbhô}.) To construct? (\textbf{\bn{নির্মাণ} \bn{করা}}, \emph{nirmān kôrā}.) A tomb? (\textbf{\bn{সমাধি}}, \emph{sômādhi}.) A dome? (\textbf{\bn{গম্বুজ}}, \emph{gômbuj}.) Architecture? (\textbf{\bn{স্থাপত্য}}, \emph{sthāpôtyô}.) A legend? (\textbf{\bn{কিংবদন্তি}}, \emph{kingbôdônti}.) A fairy tale? (\textbf{\bn{রূপকথা}}, \emph{rūpkôthā}.) A giant? (\textbf{\bn{দৈত্য}}, \emph{doityô}.) Magic? (\textbf{\bn{জাদু}}, \emph{jādu}.) To sleep? (\textbf{\bn{ঘুমানো}}, \emph{ghumāno}.) An ogre? (\textbf{\bn{রাক্ষস}}, \emph{rākkhôs}.) A fairy? (\textbf{\bn{পরী}}, \emph{pôrī}.) A witch? (\textbf{\bn{ডাইনি}}, \emph{ḍāini}.) A curse? (\textbf{\bn{অভিশাপ}}, \emph{ôbhishāp}.) A ghost? (\textbf{\bn{ভূত}}, \emph{bhūt}.)
\end{tcolorbox}
